\documentclass[11pt,a4paper]{article}
\usepackage[T1]{fontenc}
\usepackage{lmodern}
\usepackage[a4paper,margin=26mm,headheight=14pt]{geometry}
\usepackage{amsmath,amssymb,amsthm,mathtools}
\usepackage{microtype,booktabs,longtable,array,enumitem,xcolor}
\usepackage{listings,xurl,hyperref,fancyhdr}
\hypersetup{colorlinks=true,linkcolor=blue!45!black,citecolor=blue!45!black,urlcolor=blue!45!black,
 pdftitle={Local-to-global constructions in intuitionistic primitive recursive arithmetic},
 pdfauthor={Research report prepared for Vladimir Reshetnikov}}
\setlist{topsep=4pt,itemsep=2pt}
\setlength{\emergencystretch}{2em}
\newtheorem{theorem}{Theorem}[section]
\newtheorem{lemma}[theorem]{Lemma}
\newtheorem{proposition}[theorem]{Proposition}
\newtheorem{corollary}[theorem]{Corollary}
\theoremstyle{definition}
\newtheorem{definition}[theorem]{Definition}
\theoremstyle{remark}
\newtheorem{remark}[theorem]{Remark}
\newcommand{\N}{\mathbb N}
\newcommand{\Z}{\mathbb Z}
\newcommand{\Q}{\mathbb Q}
\newcommand{\R}{\mathbb R}
\newcommand{\ipra}{\mathsf{iPRA}}
\newcommand{\pra}{\mathsf{PRA}_{\mathrm{cl}}}
\newcommand{\PR}{\mathrm{PR}}
\newcommand{\code}[1]{\texttt{#1}}
\newcommand{\norm}[1]{\left\lVert#1\right\rVert}
\newcommand{\abs}[1]{\left\lvert#1\right\rvert}
\newcommand{\repo}{\url{https://github.com/VladimirReshetnikov/ProveIt}}
\pagestyle{fancy}
\fancyhf{}
\fancyhead[L]{\small Local-to-global mathematics in iPRA}
\fancyhead[R]{\small Research report}
\fancyfoot[C]{\thepage}
\lstset{basicstyle=\ttfamily\footnotesize,columns=fullflexible,breaklines=true,
 frame=single,rulecolor=\color{black!15},showstringspaces=false}

\title{Local-to-global constructions in intuitionistic primitive recursive arithmetic\\[5pt]
\large Verified iteration, polynomial differential equations, and explicit integral-point certificates}
\author{Research report prepared for Vladimir Reshetnikov}
\date{15 November 2026}

\begin{document}
\maketitle
\begin{abstract}
We extend the ProveIt collection's finite-observation approach to mathematics
in intuitionistic primitive recursive arithmetic. The inspected baseline
already contains six general framework packets, so the emphasis here is on
specific additional mechanisms and on closing a formal interface gap.
First, we give a precise interpretation discipline for finite representations
and a restricted classical-to-intuitionistic transfer theorem. Second, we
implement and kernel-check a generic primitive-recursion invariant rule whose
conclusion is an actual derivation in ProveIt's embedded \code{Prf PRAAx}
calculus. Third, we construct local solutions of rational polynomial
differential equations directly as primitive-recursive rational names, with
explicit Euler refinement bounds, uniqueness hypotheses and polynomial
conservation laws. This supplies a common analytic mechanism requested by
the classical-impossibilities packet. Fourth, we prove an explicit height
bound and complete enumeration theorem for split polynomial-branch equations,
including a uniform exact classification of a family with mixed terms.
All mathematical arguments are presented with the finite recursions and
quantifier-free invariants that support their iPRA interpretation. The large
application theorems are not yet fully formalized in the embedded calculus;
typed Lean targets and a dependency-ordered completion plan identify that
remaining work. General Siegel, Faltings and the outstanding analytic
number-theoretic estimates are not claimed.
\end{abstract}

\noindent\textbf{Reading guide.}
Section~\ref{sec:scope} compares the contribution with the inspected repository.
Section~\ref{sec:object-iteration} contains the kernel-checked object-theory rule.
The two main mathematical applications are Sections~\ref{sec:ode} and
\ref{sec:split-branches}. Section~\ref{sec:formalization} states exactly what
has been checked and what remains to be implemented.

\tableofcontents
\newpage
\section{Scope, source audit and the added mechanisms}\label{sec:scope}

The source of this investigation is the ProveIt directory
\path{Logic/PrimitiveRecursiveArithmetic}. We fixed the snapshot
\begin{center}
\code{d56cd5505fb93135b90b048bf3dcca9e73615162}.
\end{center}
Its index lists 57 research packets, including six in \code{foundations/}.
It distinguishes these delivered packets from the checked ConciseIPRA
library and says that the packets were filed without rerunning their
verification commands~\cite{repository}. That distinction governs this
report: source claims are evidence of what a packet proposes, not automatic
imports of its mathematical correctness. We inspected the foundation
packets and the particular analytic and Diophantine reports relevant below;
we have not audited all 57 packets.

The target theory is first-order intuitionistic logic with equality, the
defining equations for every fixed primitive-recursive definition, successor
injectivity and nonzeroness, and induction for quantifier-free formulas.
Parameters in induction are allowed. This is the theory represented by
\code{PRAAx} in the current ConciseIPRA source. Its existing
\code{pra\_iff\_native} theorem transfers derivations to the small native
scalar presentation~\cite{kernel,bridge}. The theorem is a transport map;
it does not create an arithmetic derivation from an ordinary Lean theorem.

\subsection{What is inherited and what is added}

The earlier packets already explain Booleanization of bounded formulas,
finite observations, explicit approximation moduli, precision propagation,
finite-window stability, simple implicit branches and several Runge-type
integer-point arguments~\cite{oldframework,embeddingpacket,finiteobservation,
evidencepacket,uniformpacket,frameworkdated}. Repeating those descriptions
would not answer the present request well. The following table makes the
additional work explicit.

\begin{longtable}{@{}>{\raggedright\arraybackslash}p{.20\linewidth}>{\raggedright\arraybackslash}p{.34\linewidth}>{\raggedright\arraybackslash}p{.40\linewidth}@{}}
\toprule
Interface & Existing coverage & Addition in this report \\
\midrule\endhead
Finite proofs & Quantifier-free induction and finite contracts; several
semantic Lean lemmas and thin object-theory adapters & A proved Lean
transformer for arbitrary primitive recursion with state, time and
parameters, returning a \code{Prf PRAAx} derivation \\
Classical reasoning & Strong emphasis on constructive presentation & A
precise criterion allowing classical first-order PRA proofs of universal
decidable conclusions, with a written negative-translation proof \\
Analysis & Fixed-function constructions, integration, implicit branches,
explicit moduli & A uniform rational-polynomial ODE constructor,
finite refinement and uniqueness bounds, and polynomial conservation \\
Integral points & Square and higher-power gaps; conditional height descent &
A separated polynomial-factor certificate for curves with mixed terms,
explicit complete enumeration, and a uniform exact example family \\
\bottomrule
\end{longtable}

These are additions to the inspected collection. We make no claim of
historical priority for negative translation, Euler's method, discrete
Gronwall estimates, conservation laws or Runge's method. The contribution is
their explicit organization for this weak first-order setting, the proofs
of the stated concrete versions, and the implemented object-theory rule.

\subsection{Four separate questions}

An embedding should answer four questions in order.
\begin{enumerate}
\item \emph{What is represented?} Specify finite codes, validity conditions,
equivalence of representatives and the exact meaning of approximations.
\item \emph{What is computed?} Give fixed PR definitions and numerical bounds,
including every search bound, precision demand and iteration count.
\item \emph{Why is it correct in the weak theory?} Supply a finite invariant
and show that its defining assertion is quantifier-free, or give its
proved Booleanization.
\item \emph{What has actually been formalized?} Distinguish a mathematical
proof, an ordinary Lean proof and a derivation in the embedded theory.
\end{enumerate}

For example, a record with a field asserting a global height bound does
not produce a height bound. A type-correct theorem target does not prove
that target. Similarly, defining a rational Euler algorithm does not prove
that its approximation errors converge. The work below supplies the latter
mathematical proofs and implements one general part of the fourth question.

\subsection{The common local-to-global pattern}

The construction has a finite state and numerical challenge. A local step
has a proved arithmetic contract. The state stores enough history or
accumulated evidence to express a quantifier-free invariant. Primitive
recursion propagates that invariant to a challenge-dependent stopping time.
Finally, a quantitative estimate converts the finite result into the desired
represented conclusion.

For an ODE, the challenge is accuracy, the state is the finite Euler table,
and the decisive estimate compares two finite meshes. For the integral-point
application, the challenge is an encoded candidate point; separated factors
force a numerical height bound, after which a finite table lists all answers.
Both use the same arithmetic proof discipline, although their mathematical
estimates are different. The framework is useful precisely because those
estimates remain visible instead of being hidden inside a universal
``certificate soundness'' assumption.

\section{Presentations, interpretations, and classical proof transfer}
\label{sec:logic}

The previous foundation packets already supply the essential numerical
tools: bounded Booleanization, finite witness search, precision-indexed
presentations, and compilation of finite observation contracts. We retain
those tools and add a precise logical interface. The distinction that matters
is between the strength of a mathematical proof and its use of classical
logic. A proof can be classical while using only the nonlogical resources of
primitive recursive arithmetic. For the conclusions specified below, its
classical logical steps can be eliminated without increasing the induction
strength. This is an application of the established negative-translation
method, not a claim of a new conservation theorem; see
\cite{moschovakis} for the broader background. We give the restricted proof
directly, because its preservation of quantifier-free induction is essential.

\subsection{The arithmetic target and its presentations}

Our target \(\ipra\) is the repository's first-order theory
\texttt{PRAAx}. Its variables range over \(\N\). Its function symbols are
finite primitive recursive definitions, represented by \texttt{PR k}; terms
are \texttt{PTm n}; formulas are \texttt{Fm PTm n}; and a completed
formalization must construct a term of type
\(\texttt{Prf PRAAx}\,\Gamma\,\varphi\).
The nonlogical axioms are defining equations, successor separation and
injectivity, and induction for quantifier-free formulas. Intuitionistic
first-order natural deduction supplies the logical rules.

An assertion about a primitive recursive approximation procedure is a
schema: its finite definition is fixed externally and then appears through
the corresponding function symbol. Its numerical parameters, dimensions,
degrees and precisions can all remain quantified. A schema is not a theorem
quantifying over a total universal evaluator for arbitrary primitive recursive
programs. The existing diagonal obstruction to such an evaluator remains in
force.

Finite lists, integers, rationals with positive denominator, polynomial
coefficient arrays and matrix entries have primitive recursive numerical
encodings. A fixed bounded formula about such data has a primitive recursive
Boolean characteristic function. Its equivalence to the characteristic
equation is proved by finite folds and quantifier-free induction, as in the
earlier framework. We use this established bounded reflection principle
throughout. It licenses a bounded invariant; it does not license every true
universal assertion about a primitive recursive function.

\begin{definition}[A numerical presentation of a sort]
\label{def:presented-sort}
A presentation of a mathematical sort consists of an arithmetic domain
predicate \(D(a)\), an arithmetic relation \(E(a,b)\) serving as equality,
and a specified code \(a_0\) with a proof of \(D(a_0)\). On \(D\), the
relation \(E\) must be provably reflexive, symmetric and transitive.
Every operation is represented by a fixed primitive recursive numerical
function, with proofs that it preserves the domains and respects \(E\).
Every additional relation is an arithmetic formula, with a proof that it
respects \(E\) in each argument.
\end{definition}

The domain and equality formulas need not be decidable. For example, for a
fixed family \(a(u,k)\) of rational observations, the assertion that \(u\)
presents a regular name and the equality of two such names are generally
universal arithmetic formulas. This is a specified family of names, not a
code space with a universal evaluator for all primitive recursive functions.
Whenever closure under an operation on name parameters is claimed, the
corresponding numerical constructor and its contracts must be supplied.
Finite-observation compilation is one way of constructing that closure.

\begin{proposition}[Interpretation through numerical presentations]
\label{prop:setoid-interpretation}
Consider a fixed finite many-sorted first-order signature with presentations
as in Definition~\ref{def:presented-sort}. Replace source equality by \(E\),
source operations and relations by their numerical interpretations, and
sorted quantifiers by
\[
 (\forall x{:}D\ A)^I=\forall x\,(D(x)\to A^I),
 \qquad
 (\exists x{:}D\ A)^I=\exists x\,(D(x)\land A^I).
\]
If the interpreted source axioms have \(\ipra\) derivations, every
intuitionistic first-order derivation from those axioms has an interpreted
\(\ipra\) derivation, with domain assumptions for its free variables.
\end{proposition}

\begin{proof}
First induct externally on terms. The closure proofs show that an interpreted
term has the appropriate domain, and the congruence proofs show that replacing
its free parameters by equivalent codes produces an equivalent result.
Then induct externally on formulas to show that the interpreted formula is
unchanged in truth under those replacements. For atomic equality this follows
from symmetry and transitivity of \(E\); for other atoms it is the supplied
relation congruence. The propositional cases follow by the intuitionistic
rules. In a quantified case keep the newly quantified code fixed, use the
induction hypothesis under its domain assumption, and introduce the
restricted quantifier.

Finally transform the given proof by induction on its finite derivation
tree. The propositional rules translate directly. Universal elimination
uses the term-domain lemma to discharge its added domain premise.
Existential introduction supplies both the term-domain proof and the
translated matrix. Existential elimination opens a witness together with
its domain proof; its eigenvariable condition is unchanged. The equality
substitution rule is the formula congruence just established. Source axioms
use their supplied interpretations. The chosen inhabitants ensure the usual
nonempty interpretation of every sort. These are external inductions on
finite syntax and proofs; they add no numerical induction schema to
\(\ipra\).
\end{proof}

This proposition gives an organization of definitions and proofs, rather
than an automatic presentation of every mathematical structure. In
particular, an induction schema in a source theory is an additional axiom
whose translation must be verified. A formula that is quantifier-free in
abstract notation may cease to be quantifier-free after interpretation.
If abstract equality means \(\forall k\,E_k(a,b)\), induction on that
equality is induction on a quantified arithmetic formula. The proposition
does not authorize it. The finite observation method must first replace it
by a numerical invariant at a fixed requested precision.

\subsection{A source theory with classical logic}

Let \(\pra\) denote the theory with exactly the same first-order language
and nonlogical axioms as \(\ipra\), but classical logic. Concretely, add to
the existing intuitionistic deduction system the axiom schema
\[
 A\lor\neg A,
 \qquad \neg A:=(A\to\bot),
\]
for arbitrary first-order formulas \(A\). Thus the source theory is not
Peano arithmetic, a higher-type theory, Lean's ambient mathematics, or set
theory. A proof in a stronger system must first be reduced to the stated
\(\pra\) before the result below can apply.

We use the elementary \(\ipra\) theorem that numerical equality is
decidable, hence stable:
\[
 (s=t)\lor\neg(s=t),
 \qquad \neg\neg(s=t)\longrightarrow s=t.
\]
The repository has an explicit arithmetic route to this fact. In
\texttt{TeqReflect.lean}, the term \texttt{Tm.eqt} reflects scalar equality
by \texttt{Q.eqt\_eq\_zero\_iff}. The cover/selector cases decide whether
the test vanishes. The interpretation \texttt{ToPRA.interp} and its round
trip transport the result to numerical \texttt{PTm} equality. This is an
object-theory proof obligation; a Lean instance of \texttt{DecidableEq}
cannot be substituted for it.

\begin{definition}[Negative translation]
\label{def:negative}
Define \(A\mapsto A^N\) externally by recursion on formulas:
\[
\begin{array}{rcl@{\qquad}rcl}
 (s=t)^N&:=&s=t,& \bot^N&:=&\bot,\\
 (A\land B)^N&:=&A^N\land B^N,
   &(A\to B)^N&:=&A^N\to B^N,\\
 (\forall x A)^N&:=&\forall x A^N,
   &(A\lor B)^N&:=&\neg(\neg A^N\land\neg B^N),\\
 (\exists x A)^N&:=&\neg\forall x\neg A^N.&&&
\end{array}
\]
This atom-preserving form uses the stability of numerical equality.
\end{definition}

Terms and numerical free variables are unchanged. Structural induction on
the formula gives
\[
 (A[\sigma])^N=A^N[\sigma]
\]
for simultaneous capture-avoiding substitution, and the corresponding
weakening identity. At a quantifier both sides use the same lifted
substitution. Moreover, \(A^N\) is quantifier-free whenever \(A\) is
quantifier-free: translating a disjunction introduces only conjunction,
implication and bottom.

\begin{lemma}[Stability of translated formulas]
\label{lem:negative-stable}
For every fixed formula \(A\), \(\ipra\) proves
\(\neg\neg A^N\to A^N\).
\end{lemma}

\begin{proof}
This is external induction on the syntax of \(A\). Equality is stable by
the arithmetic fact above. Bottom is stable by applying a hypothesis
\((\bot\to\bot)\to\bot\) to the identity implication.

Stability is preserved by conjunction: a denial of either conjunct denies
the conjunction, so double negation of the conjunction supplies double
negation of both conjuncts. Apply their stability proofs.
If \(B\) is stable, then \(C\to B\) is stable. Indeed, given
\(\neg\neg(C\to B)\) and \(C\), a denial of \(B\) would deny
\(C\to B\). We therefore obtain \(\neg\neg B\), and then \(B\).
Universal quantification preserves stability by the same argument at an
arbitrary numerical value, followed by universal introduction.

Finally every negation is stable intuitionistically. Given
\(\neg\neg\neg C\), assume \(C\); this implies \(\neg\neg C\),
contradicting the given triple negation. Both the translated disjunction
and translated existential are negations. This proves their stability
without any additional decidability assumption.
\end{proof}

\begin{theorem}[Classical-to-intuitionistic transfer at fixed strength]
\label{thm:negative-transfer}
From a \(\pra\) derivation of \(A\) with finite hypothesis context
\(\Gamma\), one can construct an \(\ipra\) derivation
\[
 \Gamma^N\vdash_{\ipra} A^N.
\]
Here \(\Gamma^N\) translates every context member. The statement requires
a source derivation; truth of the source statement alone does not suffice.
\end{theorem}

\begin{proof}
Transform the source derivation by induction on its finite proof tree.
Conjunction, implication, universal quantification and bottom use the
corresponding intuitionistic deduction rules. Substitution commutation
handles universal instantiation.

For disjunction introduction, a proof of \(A^N\) contradicts the first
conjunct of \(\neg A^N\land\neg B^N\); the other introduction is
symmetric. For disjunction elimination, the translated premises are
\[
 \neg(\neg A^N\land\neg B^N),\qquad
 A^N\to C^N,\qquad B^N\to C^N.
\]
To obtain \(C^N\), use Lemma~\ref{lem:negative-stable}. A denial of
\(C^N\) supplies both \(\neg A^N\) and \(\neg B^N\), contradicting
the first premise.

For existential introduction, an instance \(A^N(t)\) contradicts
\(\forall x\neg A^N(x)\). For existential elimination, the translated
existence premise is \(\neg\forall x\neg A^N(x)\), and the transformed
subderivation obtains \(C^N\) from \(A^N(x)\), with \(x\) fresh for
the conclusion and the other hypotheses. If \(\neg C^N\), this gives
\(\neg A^N(x)\) for arbitrary \(x\); universal introduction contradicts
the existence premise. Stability then gives \(C^N\).

The translation of excluded middle is
\[
 \neg\bigl(\neg A^N\land\neg(A^N\to\bot)\bigr).
\]
Its two prohibited conjuncts would be a formula and its negation, so this
is an intuitionistic theorem. Equality reflexivity is unchanged. For the
equality-substitution rule, all premises asserting equality of numerical
terms are unchanged, and the same rule substitutes equals in the translated
formula.

It remains to check the nonlogical axioms, which is the strength-sensitive
part of the argument. The primitive recursive defining equations,
successor separation and successor injectivity are unchanged. By the
substitution identity, the translation of the induction instance for
\(\varphi\) is precisely
\[
 \varphi^N(0)\to
 \bigl(\forall n(\varphi^N(n)\to\varphi^N(n+1))\bigr)
 \to\forall n\varphi^N(n).
\]
Since \(\varphi\) was quantifier-free, \(\varphi^N\) is also
quantifier-free. This is an allowed \(\ipra\) induction instance.
Every derivation constructor and nonlogical axiom is now covered.
\end{proof}

The transformation is syntactic. With standard finite encodings, it is a
primitive recursive operation on formula and derivation codes: each
constructor inserts a fixed proof pattern, and atomic stability instances
are substitutions into one fixed arithmetic proof. No arbitrary primitive
recursive function symbol is evaluated. A future internal theorem asserting
that this operation sends accepted \(\pra\) proofs to accepted \(\ipra\)
proofs would concern two syntactic proof checkers. It would not assert the
truth of all theorems of either system.

\subsection{Conclusions retained in their original form}

\begin{definition}[A class of retained formulas]
\label{def:retained}
Let \(\mathcal N_b\) be the class generated from fixed bounded formulas
over primitive recursive predicates by conjunction, implication and
unbounded universal quantification. We do not close this class under
unbounded existential quantification or arbitrary disjunction.
\end{definition}

\begin{lemma}[Equivalence to the translation]
\label{lem:negative-retained}
If \(A\in\mathcal N_b\), then
\(\ipra\vdash A^N\leftrightarrow A\).
\end{lemma}

\begin{proof}
First consider a bounded formula. Its bounded subformulas are decidable by
bounded Booleanization. For decidable subformulas, the negative reading of
disjunction is equivalent to ordinary disjunction. For a bounded existential,
scan its finite numerical range. If a candidate succeeds, it supplies the
existential; if all candidates fail, the resulting finite universal failure
contradicts its negative reading. This proves the equivalence for a bounded
existential without unbounded search or Markov's principle. The bounded
universal case uses pointwise equivalence. Equivalently, both readings are
equivalent to the same characteristic Boolean equation.

The remaining closure cases follow by substituting provably equivalent
subformulas in a conjunction, implication or universal quantifier. These
are external inductions on fixed formulas, not object-theory induction on
arbitrary quantified predicates.
\end{proof}

\begin{proposition}[Transfer without changing the asserted conclusion]
\label{prop:negative-conservation}
If every member of \(\Gamma\) and \(A\) belongs to \(\mathcal N_b\), then
\[
 \Gamma\vdash_{\pra} A
 \quad\Longleftrightarrow\quad
 \Gamma\vdash_{\ipra} A.
\]
\end{proposition}

\begin{proof}
Apply Theorem~\ref{thm:negative-transfer}, replace the hypotheses by their
provably equivalent translations, and then use
Lemma~\ref{lem:negative-retained} for the conclusion. The reverse direction
regards an intuitionistic proof as a classical one.
\end{proof}

In particular, a \(\pra\) proof of
\(\forall\mathbf x\,B(\mathbf x)\), with \(B\) bounded, transfers to
\(\ipra\) with the original meaning. Its intermediate proof may use
excluded middle even for unbounded formulas; the target proof contains no
new excluded-middle axiom. This applies to the logical form of exact
polynomial identities, universal finite-list bounds, arithmetic inequalities
and fixed-exponent nonexistence statements. The premise that there is a
\(\pra\) derivation remains essential: the proposition does not turn an
arbitrary published proof or an ambient Lean theorem into such a derivation.

\subsection{Witnesses recovered from bounded classical proofs}

\begin{theorem}[Bounded classical existence yields a verified selector]
\label{thm:bounded-classical-selector}
Fix a primitive recursive relation \(R(\mathbf x,y)\), a primitive
recursive numerical bound \(b(\mathbf x)\), and
\(H(\mathbf x)\in\mathcal N_b\). If \(\pra\) proves
\[
 \forall\mathbf x\left(H(\mathbf x)\to
       \exists y\le b(\mathbf x)\ R(\mathbf x,y)\right),
\]
then an explicitly defined primitive recursive function \(w\) satisfies
\[
 \ipra\vdash\forall\mathbf x\left(H(\mathbf x)\to
       w(\mathbf x)\le b(\mathbf x)\ \land\
       R(\mathbf x,w(\mathbf x))\right).
\]
\end{theorem}

\begin{proof}
The bounded existential belongs to \(\mathcal N_b\); hence the entire
displayed implication does too. Proposition~\ref{prop:negative-conservation}
first gives an \(\ipra\) proof of that implication.

Define \(w\) by scanning \(0,\ldots,b(\mathbf x)\), retaining the first
successful candidate, and returning zero if no candidate succeeds. A
separate found flag distinguishes failure from success at zero. This is
numerical primitive recursion with a given finite bound.

After testing the first \(j\) candidates, use the following bounded
invariant. Either the flag is unset and every tested candidate fails, or
the saved candidate is less than \(j\), satisfies \(R\), and every earlier
candidate fails. Bounded Booleanization makes the invariant
quantifier-free. Its base case is immediate, and its successor step is a
case split on the flag and one primitive recursive test. Quantifier-free
induction therefore proves it. The bounded existential supplied by the
transferred theorem rules out the all-failed terminal state. The saved
candidate is the required witness within its bound.
\end{proof}

This theorem enlarges the acceptable proof interface. One may prove, using
classical \(\pra\) reasoning, that a finite search space contains a desired
certificate, and then obtain the constructive certificate by the fixed
bounded selector. It is unnecessary to redesign every intermediate
classical choice in that existence proof. The substantive requirements are
the primitive recursive search bound and a proof using only the stated
nonlogical strength.

The hypotheses can still describe infinite observations. For example,
regularity of a fixed family of Cauchy names or a uniform local error
contract often has the form \(\forall\mathbf u\,C(\mathbf x,\mathbf u)=1\).
Such a formula belongs to \(\mathcal N_b\), although its validity is not a
decidable numerical test. The theorem's output and search relation remain
finite. The proofs of the existence and accuracy of an operation, and the
data required to execute it, must remain distinct.

\subsection{The boundary and the Lean completion path}

The transfer removes a logical obstacle; it supplies neither missing
quantitative estimates nor missing arithmetic induction reductions. For
example, a finiteness assertion in the form
\[
 \exists H\ \forall z\,
   (\operatorname{Point}(z)\to\operatorname{ht}(z)\le H)
\]
contains an unbounded existential. The theorem translates it to
\[
 \neg\forall H\neg\forall z\,
   (\operatorname{Point}(z)\to\operatorname{ht}(z)\le H),
\]
which does not provide \(H\). An explicit primitive recursive bound is
exactly what makes Theorem~\ref{thm:bounded-classical-selector} stronger
for its restricted conclusion. No general rule
\(\neg\neg\exists y\,R(y)\to\exists y\,R(y)\) is being assumed.
Likewise, stability of equality between represented names does not make
that equality decidable, and does not supply an apartness witness for
division or a singularity classifier for a polynomial branch.

The order of work is consequential. First interpret mathematical objects
and assertions numerically. Then replace every impermissible induction by
a verified finite invariant, fixing observation precision where necessary.
After this arithmetic audit, Theorem~\ref{thm:negative-transfer} can remove
classical logic from a \(\pra\) derivation. Finally a bounded existential
can be realized by Theorem~\ref{thm:bounded-classical-selector}.
Performing negative translation before the arithmetic audit does not cure
induction on an unbounded observation formula.

The companion \texttt{NegativeSyntax.lean} completes the syntax and
axiom-translation layer. It compiles with Lean 4.32.0; its checked declarations
use neither \texttt{sorryAx} nor \texttt{Classical.choice}. The audit records
Lean's standard \texttt{propext} and \texttt{Quot.sound} dependencies in the
metatheory. The completed and remaining steps are concrete:
\begin{enumerate}
\item The file defines \texttt{negative} on the existing \texttt{Fm PTm n},
  and \texttt{ClassicalPRAAx} with one constructor for a \texttt{PRAAx}
  instance and one for excluded middle. It proves the syntax identities
  \texttt{negative\_subst}, \texttt{negative\_shift} and
  \texttt{negative\_qf} by structural induction.
\item Produce the object proof of numerical equality stability through the
  arithmetic route described above. Then construct an actual
  \texttt{Prf PRAAx} proof of stability for every translated formula, using
  the proof patterns in Lemma~\ref{lem:negative-stable}.
\item The file also translates every arithmetic and excluded-middle axiom
  constructor into an actual \texttt{Prf PRAAx} derivation. For induction,
  it invokes \texttt{PRAAx.ind} on \texttt{negative phi};
  \texttt{negative\_qf} certifies the precise side condition. The excluded
  middle translation is proved directly by intuitionistic contradiction.
\item The remaining main task is to prove by induction on the existing
  \texttt{Prf} derivation the target
\[
 \texttt{Prf ClassicalPRAAx}\ \Gamma\ \varphi
 \ \longrightarrow\
 \texttt{Prf PRAAx}\ (\Gamma\texttt{.map negative})\
       (\texttt{negative}\ \varphi).
\]
  The \texttt{eqSubst} case uses the same constructor on the translated
  formula; the quantifier cases use the substitution and weakening
  identities. The critical disjunction and existential elimination cases
  are exactly those written out in Theorem~\ref{thm:negative-transfer}.
\item The file defines retained formulas with quantifier-free leaves.
  Their equivalence and transfer theorems remain named goal propositions,
  as does the full proof translator. The bounded-leaf extension then uses a
  separately proved Booleanization adapter. Derive transfer of retained
  sequents using
  \texttt{Prf.trans} and the equivalence proofs.
\item Implement first-success search using \texttt{PR} symbols and prove its
  bounded found/no-success invariant. The object recurrence rule developed
  in this report supplies the corresponding induction step. This produces
  the selector theorem at the actual \texttt{Prf PRAAx} level.
\end{enumerate}

The human proof of the logical transfer and the selector theorem is complete
above. Their full Lean proof translator and equality-stability adapter are
specified obligations, not completed formalizations. The independently
compiled object recurrence result reported later should not be confused
with completion of this entire translation pipeline.

\section{A verified object-theory rule for primitive recursion}
\label{sec:object-iteration}

This section supplies a concrete formal advance at the common endpoint of
the proof packets. The theorem is an inference rule about embedded iPRA
derivations. It does not appeal to the soundness or completeness of iPRA in
the standard natural numbers.

\subsection{Mathematical statement and proof}

Fix a tuple of numerical parameters $\vec x$ of length $k$, a PR initial
state function $f(\vec x)$ and a PR update function $g(t,s,\vec x)$. Define
\begin{align*}
 u(0,\vec x)&=f(\vec x),\\
 u(t+1,\vec x)&=g(t,u(t,\vec x),\vec x).
\end{align*}
The state $s$ is one number; tuples, finite tables and certificates may be
encoded in that number. Let $P(s,t,\vec x)$ be a quantifier-free formula.
Let $\Gamma(\vec x)$ be any finite list of assumptions, possibly with
unbounded quantifiers. They are fixed parameters in the induction and do
not become part of its induction formula.

\begin{theorem}[Primitive-recursion invariant rule]\label{thm:object-rule}
Suppose there are iPRA derivations of
\begin{align}
 \Gamma(\vec x)&\vdash P(f(\vec x),0,\vec x),\label{eq:object-base}\\
 P(s,t,\vec x),\Gamma(\vec x)&\vdash
 P(g(t,s,\vec x),t+1,\vec x).\label{eq:object-step}
\end{align}
Then there is an iPRA derivation of
\[
 \Gamma(\vec x)\vdash\forall t\,P(u(t,\vec x),t,\vec x).
\]
Moreover the derivation is obtained uniformly from the two supplied
derivations, the finite PR syntax for $f,g$, and the formula $P$.
For any PR term $b(\vec x)$ it gives
\[
 \Gamma(\vec x)\vdash
 P(u(b(\vec x),\vec x),b(\vec x),\vec x).
\]
\end{theorem}

\begin{proof}
Put $I(t,\vec x)=P(u(t,\vec x),t,\vec x)$. Substitution of terms in a
quantifier-free formula preserves quantifier-freeness, so $I$ is an allowed
induction formula. At zero, the defining equation of primitive recursion
and equality substitution turn \eqref{eq:object-base} into $I(0,\vec x)$.
For the induction step, instantiate the freely available state variable
$s$ in \eqref{eq:object-step} by $u(t,\vec x)$. Its premise becomes
$I(t,\vec x)$ and its conclusion becomes
$P(g(t,u(t,\vec x),\vec x),t+1,\vec x)$. The successor defining equation
of primitive recursion rewrites this as $I(t+1,\vec x)$.
Quantifier-free induction, with the unchanged assumptions $\Gamma$, gives
$\forall t\,I(t,\vec x)$. Universal elimination at $b(\vec x)$ proves
the final assertion. All proof transformations are syntactic substitution,
weakening, equality substitution and one instance of the permitted induction
rule. No induction on a formula containing a quantifier has been used.
\end{proof}

The important point is that \eqref{eq:object-step} is proved once with
arbitrary $s,t,\vec x$. A separately supplied proof for each standard
numeral $t$ is a different input and would not justify the theorem. Likewise,
the bounds and challenges in $\vec x$ remain variables in a single theorem.

\subsection{The exact Lean implementation}

The file \code{lean/ObjectKernel.lean} contains the current repository
declarations from \code{Syntax.lean}, \code{Deduction.lean} and
\code{PR.lean}, with imports replaced by \code{Std} and the notation
\code{Nat} supplied locally. Their declaration bodies are unchanged; hashes
and the inherited license are supplied. This avoids requiring the full
Mathlib dependency closure to check the new rule, while retaining the exact
embedded theory's syntax, deduction constructors and axioms.

\code{ObjectIteration.lean} fixes the following variable convention:
\[
 P:\code{Fm PTm}\,(k+2),\qquad
 0=\text{state},\quad1=\text{time},\quad2,3,\ldots=\text{parameters}.
\]
In the PR step symbol, the order is instead the repository's required
$(\text{time},\text{state},\text{parameters})$. The file explicitly performs
that permutation. This is a substantive scoping detail: exchanging those
arguments accidentally could produce a plausible-looking but incorrect
compiler interface.

The key declarations, using shortened qualification, are:
\begin{lstlisting}
runTerm (f : PR k) (g : PR (k + 2))
  (t : PTm n) (xs : Fin k -> PTm n) : PTm n

instStateTime (P : Fm PTm (k + 2))
  (s t : PTm n) (xs : Fin k -> PTm n) : Fm PTm n

primitiveRecursionInvariant
  (f : PR k) (g : PR (k + 2))
  (P : Fm PTm (k + 2)) (hP : P.QF)
  (Gamma : List (Fm PTm k))
  (h0 : Prf PRAAx Gamma (baseFormula f P))
  (hS : Prf PRAAx
    (P :: (Gamma.map Fm.shift).map Fm.shift)
    (nextFormula g P)) :
  Prf PRAAx Gamma (Fm.all (runInvariant f g P))
\end{lstlisting}
The displayed code is a readable signature; the file uses the repository's
sequent notation and contains the complete proof. Its proof first establishes
substitution identities for state and time, then uses
\code{PRA.precZero}, \code{PRA.precSucc}, \code{Prf.substWeaken},
\code{Prf.eqSubst} and \code{PRA.qfInd}. The separate theorem
\code{primitiveRecursionAt} instantiates the fuel term.

\subsection{A closed test of the proof rule}

\code{ObjectExample.lean} chooses $f=0$, the update $g(t,s)=s+1$ and
$P(s,t)\equiv s=t$. Initialization is reflexivity. The local step follows
from the explicit PR composition and projection equations and congruence
of successor. The generic rule then constructs the closed object theorem
\[
 \ipra\vdash\forall t\,[u(t)=t].
\]
This small example checks the argument ordering, free-variable substitution
and genuine induction endpoint. It is not offered as a new arithmetic
result. The more useful artifact is the general derived rule itself.

\subsection{What the result does and does not automate}

The rule accepts \emph{proofs} of initialization and preservation. It does
not accept semantic truth values in their place. Once the integer, rational
and array operations are represented by PR syntax and their local lemmas
are proved, a large bounded algorithm can use the rule repeatedly without
rebuilding its induction plumbing. Bounded quantifiers in an invariant must
first be replaced by their Boolean tests with a proved equivalence.

This is a verified proof transformer, not a completed compiler from Python,
arbitrary Lean, or informal mathematics. In particular, the ODE and
integral-point estimates below still require their own arithmetic
object-theory proofs. The rule resolves their common iteration interface;
it does not manufacture those estimates.

The standalone build checks the new rule against the exact arithmetic
kernel declarations. Integration into ProveIt replaces the standalone
kernel import with \code{import ConciseIPRA.PR}. After that build,
\code{pra\_iff\_native.mp} transports the resulting sequent. The optional
repository adapter is supplied as source and is explicitly excluded from
the standalone verification claim.

\section{Polynomial differential equations from finite arithmetic}
\label{sec:ode}

The three-impossibilities packet asks whether the analytic infrastructure
for the Hermite identity and disk-area calculation can be organized around
a generic theorem for polynomial differential equations with explicit
majorants \cite{impossibilitiespacket}.  We give such a construction, together
with a finite algebraic method for proving its conserved quantities.  This
extends the earlier libraries of names, finite integration, and implicit
branches \cite{oldframework}.  The new construction produces a primitive
recursive approximation scheme and its finite error estimates; its
existence proof never invokes completeness of the real numbers.

\subsection{Finite polynomial data and a constructed local horizon}

A rational polynomial vector field in dimension $d\ge1$ is a finite table
\[
 F_i(x)=\sum_{\alpha}c_{i,\alpha}x^\alpha,
 \qquad x^\alpha=\prod_{j<d}x_j^{\alpha_j},\qquad
 |\alpha|=\sum_{j<d}\alpha_j.
\]
Rational arithmetic, exponent vectors, and all the tables are numerically
coded.  Polynomial evaluation is uniform primitive recursion on this
particular finite syntax; it is not evaluation of arbitrary primitive
recursive program codes.  Dimension, degree, coefficient sizes, and
precision may all vary.  We use the norm
$\lVert x\rVert=\max_{i<d}|x_i|$.

For rational $B\ge1$, compute
\begin{align}
 M&=\max\left(1,\max_{i<d}\sum_\alpha
                  |c_{i,\alpha}|B^{|\alpha|}\right),\label{eq:ode-M}\\
 L&=\max_{i<d}\sum_{|\alpha|\ge1}
                  |c_{i,\alpha}|\,|\alpha|B^{|\alpha|-1}.
                  \label{eq:ode-L}
\end{align}

\begin{lemma}[Finite polynomial bounds]
\label{lem:ode-polybound}
If $\lVert x\rVert,\lVert y\rVert\le B$, then
\[
 \lVert F(x)\rVert\le M,\qquad
 \lVert F(x)-F(y)\rVert\le L\lVert x-y\rVert.
\]
\end{lemma}
\begin{proof}
The first assertion is the triangle inequality applied to each monomial.
For the second, write a monomial as a product of $|\alpha|$ factors, with
coordinates repeated as prescribed by $\alpha$.  Replace the factors at
$x$ by the factors at $y$ one at a time.  The resulting telescoping identity
has $|\alpha|$ terms, each of absolute value at most
$B^{|\alpha|-1}\lVert x-y\rVert$.  Multiplication by coefficient absolute
values and finite summation give the result.  These are finite identities
and inequalities, including when dimension and degree vary.
\end{proof}

Given any rational initial vector $x_0$, the choices
\begin{equation}
 B=1+\lVert x_0\rVert,\qquad M,L\text{ from
 \eqref{eq:ode-M}--\eqref{eq:ode-L}},\qquad T=\frac1{2M}
 \label{eq:ode-localhorizon}
\end{equation}
construct a positive rational horizon.  For $N=2^n$ and $h=T/N$, define
\begin{equation}
 x_{n,0}=x_0,\qquad x_{n,j+1}=x_{n,j}+hF(x_{n,j})\quad(j<N).
 \label{eq:ode-euler}
\end{equation}
Induction gives
\begin{equation}
 \lVert x_{n,j}-x_0\rVert\le jhM\le TM=\tfrac12\quad(j\le N).
 \label{eq:ode-invariant}
\end{equation}
At each step the preceding estimate puts the current point in the box,
where Lemma~\ref{lem:ode-polybound} applies.  Thus both the local horizon
and its validity proof are produced from the finite input.  They are not
extra existence hypotheses.

The subsequent theorems also work on any larger rational horizon $T>0$
for which a uniform arithmetic proof establishes that the Euler nodes
remain in $\lVert x\rVert\le B$.  The finite inequality
$\lVert x_0\rVert+TM\le B$ is one sufficient certificate.  Concrete linear
and triangular systems below admit sharper invariant proofs.  Restricting
to meshes $2^{n+b}$, for a supplied natural $b$, is harmless: every bound
below remains true with $2^{-n}$ as an overestimate of the mesh factor.

\subsection{Discrete Gronwall with an entirely rational growth bound}

\begin{lemma}[Rational growth bound]
\label{lem:ode-growth}
For $a\in\Q$, $a\ge0$, and integer $k\ge1$,
\[
 \left(1+\frac ak\right)^k\le3^{\lceil a\rceil}.
\]
\end{lemma}
\begin{proof}
Let $A=\lceil a\rceil$.  Finite Bernoulli gives
$1+a/k\le1+A/k\le(1+1/k)^A$.  The binomial expansion gives
\[
 (1+1/k)^k\le\sum_{j=0}^k\frac1{j!}\le3,
\]
because $\binom{k}{j}\le k^j/j!$ and $j!\ge2^{j-1}$ for $j\ge2$;
the last step is a finite geometric sum.  Raise the Bernoulli inequality
to the $k$th power.  If $A=0$, then $a=0$ and the claim is immediate.
\end{proof}

\begin{lemma}[Finite Gronwall estimate]
\label{lem:ode-gronwall}
Suppose $e_0,\ldots,e_k$ and $\delta_0,\ldots,\delta_{k-1}$ are
nonnegative rationals and
$e_{j+1}\le r e_j+\delta_j$, where $r\ge1$.  Then
\[
 e_j\le r^j\left(e_0+\sum_{i<j}\delta_i\right)\quad(j\le k).
\]
In particular, if $r=1+a/k$, $k\ge1$, and $\delta_i\le\delta$, then
\begin{equation}
 e_j\le3^{\lceil a\rceil}(e_0+k\delta)\quad(j\le k).
 \label{eq:ode-gronwall}
\end{equation}
\end{lemma}
\begin{proof}
The first inequality is a finite induction.  In its step, multiply the
induction bound by $r$ and use
$\delta_j\le r^{j+1}\delta_j$.  For the second assertion, use
$r^j\le r^k$, Lemma~\ref{lem:ode-growth}, and the finite sum bound.
After rational coding and Booleanization of the bounded input conditions,
the invariant is quantifier-free.  It involves no analytic assertion.
\end{proof}

\subsection{A primitive recursive family of solution names}

Interpolate the nodes of \eqref{eq:ode-euler} by
\begin{equation}
 U_n(t)=x_{n,j}+(t-jh)F(x_{n,j})
 \quad(jh\le t\le(j+1)h),\qquad U_n(T)=x_{n,N}.
 \label{eq:ode-polygon}
\end{equation}
The containing cell is found by integer division.  At grid boundaries
the two formulas agree.  The invariant box is convex, so it contains the
whole polygon.  Every polygon is $M$-Lipschitz: split any rational time
interval at its finitely many mesh points and sum the bounds on the
constant slopes.

Compute
\begin{equation}
 A=\lceil LT\rceil,\qquad G=3^A,
 \qquad K=\max(1,\lceil GLMT^2\rceil).
 \label{eq:ode-constants}
\end{equation}

\begin{theorem}[Uniform finite refinement estimate]
\label{thm:ode-refine}
For every rational $0\le t\le T$,
\begin{equation}
 \lVert U_{n+1}(t)-U_n(t)\rVert
       \le\frac{GLMT^2}{2\,2^n}.
 \label{eq:ode-refine}
\end{equation}
Consequently, for all $m\ge n$,
\begin{equation}
 \lVert U_m(t)-U_n(t)\rVert\le K2^{-n}.
 \label{eq:ode-cauchy}
\end{equation}
\end{theorem}
\begin{proof}
Put $N=2^n$, $h=T/N$, and compare the two polygons at times $kh/2$,
$k\le2N$.  Write $y_k=U_n(kh/2)$, $z_k=U_{n+1}(kh/2)$ and
$e_k=\lVert y_k-z_k\rVert$.  At an even $k$ the coarse slope is
$F(y_k)$.  At an odd $k$, the point $y_k$ is a half-step from a coarse
node $x$, so $\lVert y_k-x\rVert\le Mh/2$.  Thus the coarse slope
differs from $F(y_k)$ by at most $LMh/2$.  In both cases,
\[
 e_{k+1}\le(1+Lh/2)e_k+LMh^2/4,
 \qquad e_0=0.
\]
Apply \eqref{eq:ode-gronwall} with $2N$ steps and $a=LT$.  It gives
$e_k\le GLMT^2/(2N)$ at every fine node.  On a fine interval both
polygons are affine, and their difference is the convex interpolation
of its endpoint values.  The same norm bound therefore holds throughout
the interval.  This proves \eqref{eq:ode-refine}.

For \eqref{eq:ode-cauchy}, telescope the finite list
$U_n(t),\ldots,U_m(t)$ and use
$\sum_{j=n}^{m-1}2^{-j-1}\le2^{-n}$.  No infinite sum or completeness
principle occurs.
\end{proof}

For integer $K\ge1$, let $\ell(K)$ be the least $r\le K$ such that
$K\le2^r$, found by bounded search.  With
\begin{equation}
 c=\ell(K)+3,\qquad X(p,t)=U_{p+c}(t),
 \label{eq:ode-name}
\end{equation}
the estimate \eqref{eq:ode-cauchy} gives
\[
 \lVert X(q,t)-X(p,t)\rVert\le2^{-p-3}\qquad(q\ge p).
\]
This is a uniform regular name, with spare precision.  The mathematical
object produced by the theorem is this primitive recursive rational
evaluator and these contracts, not an independently postulated real limit.

For completeness, the time input can itself be represented.  Given a fixed
regular time-name evaluator $\tau(p)$ representing a point in $[0,T]$,
clip its rational approximations to $[0,T]$.  Clipping does
not increase distance.  Put
$d_0=\ell(\max(1,\lceil M\rceil))+3$ and use
\[
 U_{p+c}\bigl(\operatorname{clip}(\tau(p+d_0))\bigr).
\]
Split the difference of two outputs into a same-time mesh error and a
same-mesh time error.  The former is at most $2^{-p-3}$, and the latter,
by $M$-Lipschitz continuity, is at most $2^{-p-3}+2^{-q-3}$ when $q\ge p$.
This gives the regular-name contract.  The evaluator for $\tau$ is fixed,
or supplied through a finite table with explicit demands.

\subsection{The solution contract and its derivative modulus}

The following definition states precisely which differential-equation
claim is proved.  In particular, it keeps the uniform error data visible.

\begin{definition}[Controlled solution name]
\label{def:controlled-flow}
For fixed finite data $F,x_0,B,T$ and rational $C\ge0$, a controlled
solution name is a primitive recursive rational evaluator $Y(p,t)$,
for rational $0\le t\le T$, such that:
\begin{enumerate}
\item it is regular in $p$, uniformly in rational $t$: for all $p,q$,
      $\lVert Y(p,t)-Y(q,t)\rVert\le2^{-p}+2^{-q}$;
\item $\lVert Y(p,t)\rVert\le B$ and
      $\lVert Y(p,0)-x_0\rVert\le2^{-p}$;
\item for rational $q,r\in[0,T]$,
\begin{equation}
 \begin{split}
 &\lVert Y(p,r)-Y(p,q)-(r-q)F(Y(p,q))\rVert\\
 &\hspace{30mm}\le C|r-q|^2+2^{-p}.
 \end{split}
 \label{eq:ode-controlled}
\end{equation}
\end{enumerate}
For a fixed evaluator these are first-order arithmetic formulas.
The construction of its code and the proof of its contracts are separate
from the assertion that it is a solution name.
\end{definition}

\begin{theorem}[Constructed controlled flow]
\label{thm:ode-existence}
The evaluator \eqref{eq:ode-name} is a controlled solution name with
$C=LM/2$.  Hence every rational polynomial vector field and rational
initial vector have a uniformly constructed positive local horizon and
a primitive recursive controlled solution there.
\end{theorem}
\begin{proof}
Regularity, the box bound and the initial value have already been proved.
Let $v=|r-q|$.  On the cell containing $s$, its left endpoint is within
$Mh$ of $U_n(s)$.  The latter is within $M|s-q|$ of $U_n(q)$.
Lemma~\ref{lem:ode-polybound} bounds the difference between the polygon
slope and $F(U_n(q))$ by $LM(h+|s-q|)$.

Split the interval between $q$ and $r$ at its finitely many mesh points.
On each resulting subinterval the slope difference is constant.  Evaluate
the affine upper bound at that subinterval's rational midpoint and multiply
by its length.  Summing these inequalities and the constant polygon slopes
yields
\begin{equation}
 \begin{split}
 &\lVert U_n(r)-U_n(q)-(r-q)F(U_n(q))\rVert\\
 &\hspace{30mm}\le LM\bigl(v^2/2+hv\bigr).
 \end{split}
 \label{eq:ode-localdefect}
\end{equation}
For $q\le r$, the midpoint-times-length sum of $s-q$ is exactly
$(r-q)^2/2$, by the difference-of-squares identity on each subinterval.
Thus this uses only finite rational arithmetic, with no general positivity
theorem for integrals.  Reversing the orientation proves
the same estimate for $r<q$.

Substitute $n=p+c$, use $v\le T$, and observe
$LMT^2\le K$ and $K2^{-c}\le1/8$.  The extra mesh term is at most
$2^{-p-3}$, so \eqref{eq:ode-controlled} follows with $C=LM/2$.
Finally \eqref{eq:ode-localhorizon} supplies all hypotheses for arbitrary
finite rational polynomial input.
\end{proof}

The name-level interpretation of \eqref{eq:ode-controlled} is
\[
 \lVert Y(r)-Y(q)-(r-q)F(Y(q))\rVert\le C|r-q|^2.
\]
For nonzero rational $r-q$, the difference-quotient error is at most
$C|r-q|$.  Thus
$|r-q|\le2^{-k}/\max(1,\lceil C\rceil)$ is a derivative modulus,
with one-sided derivatives at endpoints.  This interpretation can be
proved entirely as all-precision rational tests using
\eqref{eq:ode-controlled}; no exact real sign or equality decision is used.
The same contract gives a time-continuity modulus, since for $|r-q|\le1$
the increment is bounded by $(M+C)|r-q|+2^{-p}$ at precision $p$.

\subsection{Uniqueness proved at one finite precision}

\begin{theorem}[Quantitative uniqueness]
\label{thm:ode-unique}
Two controlled solution names for the same $F,x_0,B,T$, with respective
defect constants $C_1,C_2$, represent the same function on $[0,T]$.
There is an explicit primitive recursive schedule for the mesh and name
precision used to verify equality at any requested accuracy.
\end{theorem}
\begin{proof}
Fix a rational $t\in[0,T]$, a positive integer $N$, and a name precision
$p$.  Compare the rational approximations to the two candidates at
$jt/N$, $j\le N$.  If $e_j$ is their distance, the two instances of
\eqref{eq:ode-controlled} give
\[
 e_{j+1}\le(1+Lt/N)e_j+(C_1+C_2)(t/N)^2+2\,2^{-p},
 \qquad e_0\le2\,2^{-p}.
\]
The finite Gronwall estimate therefore gives
\begin{equation}
 e_N\le G\left((C_1+C_2)T^2/N+2(N+1)2^{-p}\right).
 \label{eq:ode-uniquefinite}
\end{equation}
For a requested accuracy $2^{-k}$ choose a dyadic $N$ at least
$4G(C_1+C_2)T^2 2^k$ (and at least one), then choose $p$ with
$2^p\ge16G(N+1)(1+LT)2^k$.  Both choices use bounded logarithms of explicit
integers.  Each term in \eqref{eq:ode-uniquefinite} is at most
$2^{-k-2}$.  The resulting arbitrarily accurate agreement, together with
regularity, is exactly equality of the two name families: a primitive
recursive common reindexing makes their rational outputs differ by at
most $2^{-k}$ for every $k$.  Regularity transfers this agreement to every
finer requested observation.

In the object-theory proof, $N,p,t,k$ are numerical parameters and the
induction is on the rational error table $e_j$.  It is not induction on
the assertion that two real functions are equal.
\end{proof}

The theorem deliberately specifies controlled candidates.  It does not
extract a primitive recursive modulus from qualitative differentiability
of an arbitrary external function.  Its useful feature is that the
constructed solution supplies the requisite modulus automatically.

A related finite comparison is useful when composing constructions.  If
two polynomial fields have a common Lipschitz bound $L$ and a coefficient
estimate proves $\lVert F(x)-\widetilde F(x)\rVert\le\varepsilon$ on the
box, their Euler trajectories with initial separation $\eta$ satisfy
\[
 e_{j+1}\le(1+Lh)e_j+h\varepsilon.
\]
Their represented flows therefore obey
\begin{equation}
 \lVert X_F(t)-X_{\widetilde F}(t)\rVert
       \le G(\eta+T\varepsilon).
 \label{eq:ode-perturbation}
\end{equation}
An additional step error $\delta_j$ adds
$G\sum_j\delta_j$ to the bound.  This gives a predetermined precision
budget for rounded steps or Cauchy-named coefficients, provided the
common domain and invariance contracts are supplied.  The unconditional
local theorem above already covers arbitrary finite rational polynomial
codes.

\subsection{Finite algebraic identities produce conserved quantities}

Let $H(x)=\sum_\alpha b_\alpha x^\alpha$ be a rational polynomial.
Its formal partial derivatives are finite coefficient tables.

\begin{theorem}[Polynomial conservation compiler]
\label{thm:ode-conservation}
If the coefficientwise polynomial identity
\begin{equation}
 \sum_{i<d}(\partial_iH)(x)F_i(x)=0
 \label{eq:ode-conservationcert}
\end{equation}
holds, then the constructed flow satisfies
$H(X(t))=H(x_0)$ for every represented time $t\in[0,T]$.
\end{theorem}
\begin{proof}
Put $D=\max(1,B,MT)$ and compute the rational
\[
 R_H=\sum_{|\alpha|\ge2}|b_\alpha|\,2^{|\alpha|}
                                  D^{|\alpha|-2}.
\]
The multivariate binomial expansion gives the finite remainder estimate
\begin{equation}
 |H(x+v)-H(x)-\nabla H(x)\cdot v|\le R_H\lVert v\rVert^2
 \label{eq:ode-polyremainder}
\end{equation}
when $\lVert x\rVert\le B$ and $\lVert v\rVert\le MT$.
Indeed, degree-zero and degree-one terms in $v$ are precisely the
subtracted terms.  In every remaining term retain two factors of $v$
and bound the other factors by $D$.  For a monomial of total degree $e$,
the sum of the relevant binomial coefficients is at most $2^e$.
Finite summation proves \eqref{eq:ode-polyremainder}.

Take $v=hF(x)$.  Identity \eqref{eq:ode-conservationcert} makes the
linear term zero, so a full Euler step changes $H$ by at most
$R_HM^2h^2$.  At time $t=jh+\theta$, $0\le\theta\le h$, telescope
the $j$ full steps and the final partial step.  Since
$jh^2+\theta^2\le h(jh+\theta)\le hT$,
\begin{equation}
 |H(U_n(t))-H(x_0)|\le R_HM^2T^2\,2^{-n}.
 \label{eq:ode-conservationerror}
\end{equation}
The polynomial $H$ has an explicit Lipschitz bound on the box, by the
proof of Lemma~\ref{lem:ode-polybound}.  Combining it with
\eqref{eq:ode-cauchy} and \eqref{eq:ode-conservationerror}, at an
arbitrary specified precision, proves equality of the represented values.
The represented-time extension uses the same continuity bound.
\end{proof}

The certificate \eqref{eq:ode-conservationcert} is checked by finite
polynomial multiplication and rational coefficient comparison.  The
passage from that certificate to an exact invariant uses only the
displayed finite error bounds.  This is a useful general separation:
algebra finds the identity, while one common arithmetic theorem transports
it to the represented flow.

If instead $\nabla H\cdot F=J$, retaining the linear term gives the
corresponding integral identity.  One can also add an accumulator coordinate
$z'=J(x)$ and apply the zero-derivative theorem to $H(x)-z$.
The accumulator's increments are finite Riemann sums with an explicitly
vanishing defect.  Thus its identification with the integral uses the
existing finite-sum integration interface, without invoking an existence
theorem for integrals of arbitrary functions.

\subsection{A common engine for the Hermite identity and circle calculus}

\paragraph{Exponential and oscillator names.}
For $x'=x$, $x(0)=1$, on any rational horizon $T>0$, finite Euler states
are bounded by
$(1+T/N)^j\le3^{\lceil T\rceil}$.  Thus a uniform invariant box is
available directly from Lemma~\ref{lem:ode-growth}, without imposing
the smaller local horizon.  At $T=1$ the endpoint approximants are
$(1+2^{-n})^{2^n}$.  The generic construction supplies a name, its
derivative contract, and its uniqueness property.  Existing exponential
names with the same controlled residual and initial value can be identified
by Theorem~\ref{thm:ode-unique}.

For the oscillator
\[
 c'=-s,\qquad s'=c,\qquad(c(0),s(0))=(1,0),
\]
the sup norm of a step grows by at most $1+h$, so the same invariant
box works.  The coefficient identity
\[
 \partial_c(c^2+s^2)(-s)+\partial_s(c^2+s^2)c
       =-2cs+2sc=0
\]
gives $c(t)^2+s(t)^2=1$ through Theorem~\ref{thm:ode-conservation}.
Once the standard sine and cosine names satisfy the controlled derivative
contract, one uniqueness theorem identifies both constructions.

\paragraph{The finite Hermite antiderivative identity.}
Let $f\in\Q[t]$ have degree at most $m$, and compute
\[
 Q=\sum_{j=0}^{m}f^{(j)}.
\]
Formal differentiation is a finite coefficient operation, and telescoping
gives $Q-Q'=f$.  Consider the polynomial system
\begin{equation}
 \tau'=1,\qquad E'=-E,\qquad J'=E f(\tau),
 \qquad (\tau(0),E(0),J(0))=(0,1,0).
 \label{eq:hermite-flow}
\end{equation}
For any supplied rational horizon $T>0$, begin with meshes satisfying
$h\le1$, accomplished by a bounded logarithm.  The Euler invariant is
explicit: $0\le\tau_j\le T$, $0\le E_j\le1$, and
\[
 |J_j|\le jh F_0\le TF_0,\qquad
 F_0=\sum_{k\le m}|f_k|\max(1,T)^k.
\]
The bound on $E_j$ follows from $E_{j+1}=(1-h)E_j$; the bound on $J_j$
is a finite induction.  Therefore the generic flow theorem applies on
the entire supplied horizon.

The polynomial $H(\tau,E,J)=E Q(\tau)+J$ satisfies
\[
 \nabla H\cdot(1,-E,Ef(\tau))
   =E\bigl(Q'(\tau)-Q(\tau)+f(\tau)\bigr)=0.
\]
The conservation theorem proves
\begin{equation}
 J(t)=Q(0)-E(t)Q(t).
 \label{eq:hermite-flowidentity}
\end{equation}
Here $\tau(t)=t$, already exactly true of the Euler polygons, and $E$
is the controlled solution of $E'=-E$, $E(0)=1$.  The accumulator $J$
is the integral name of $E(t)f(t)$: sum its controlled increments over
a rational mesh.  At finite table precision $p$, the total defect is
bounded by $C_JT^2/N+(N+1)2^{-p}$, where $C_J$ is its controlled-increment
constant.  First choose $N$ to make the mesh term small, then choose $p$
to make the second term small.  The product-evaluation error has an
explicit coefficient bound and is included by increasing $p$ further.
Thus the sums converge with an explicit modulus to $J(t)$.
Identifying $E$ with the existing $e^{-t}$ name by uniqueness gives the
usual Hermite identity
\begin{equation}
 \int_0^t f(u)e^{-u}\,du=Q(0)-e^{-t}Q(t).
 \label{eq:hermite-identity}
\end{equation}
The integral notation in \eqref{eq:hermite-identity} denotes this
explicit finite-sum name.  Formula \eqref{eq:hermite-flowidentity} was
proved first, using only finite polynomial identities and Euler estimates.
The divisibility arguments and the final transcendence contradiction in
the original application remain separate arithmetic steps.

\paragraph{Circle-area accumulator.}
Add coordinates $\tau,A$ to the oscillator with
$\tau'=1$ and $A'=(c^2+s^2)/2$, initially zero.  The same finite
growth bound supplies an Euler box for $c,s$, and
$|A_j|\le T\,3^{2\lceil T\rceil}$ is a sufficient accumulator bound.
The two polynomial certificates are
\[
 \nabla(c^2+s^2)\cdot F=0,
 \qquad
 \nabla\bigl(A-\tau(c^2+s^2)/2\bigr)\cdot F=0.
\]
Consequently $c^2+s^2=1$ and $A(t)=t/2$.  Since
$(c s'-s c')/2=(c^2+s^2)/2$, $A$ is the signed sector-integral
accumulator.  The identification of a full period with $2\pi$ and the
geometric interpretation of that integral as disk area are still supplied
by the relevant geometric argument.  What has been resolved is the
common analytic construction and its exact polynomial identities, not
every step of those applications.

\subsection{The iPRA induction audit}

All algorithms here are primitive recursive uniformly in their finite
inputs.  The large value $2^{p+c}$ is harmless for primitive recursiveness;
the theorem makes no feasibility claim.  The required inductions have
explicit finite invariants: the Euler box estimate, the Gronwall bound
with partial error sum, the paired fine/coarse-node comparison, the
geometric telescoping bound, and the accumulated conservation defect.
Rational inequalities are reduced to integer inequalities by positive
denominators.  Finite coordinate, monomial, and mesh quantifiers are
Booleanized before induction.

For equality or differentiability of names, first fix a numerical
precision challenge.  The displayed bounds compute a sufficient mesh
and finite evaluation table.  Prove the resulting finite arithmetic
assertion, and only then introduce the challenge universally.  This is
the point at which one avoids induction on an analytic formula.

The mathematical construction therefore has a concrete path to
$\mathsf{Prf}\,\mathsf{PRAAx}$ proofs: code the finite algorithms, prove
their defining equations, and implement those bounded invariants in the
existing object calculus.  The complete deep Lean derivations are not
claimed here.  Ambient Lean correctness of a rational Euler algorithm
and an object-language iPRA proof remain distinct deliverables.

\section{An algebraic producer of integral-point bounds}
\label{sec:split-branches}

A framework becomes more useful when some of its global quantitative
inputs can themselves be produced from finite algebra.  The earlier
Siegel packet proves a square-completion certificate
\(D^2 f=A^2+R\), and the earlier general framework already extends it
to all powers \(D^m f=A^m+R\) with a degree gap
\cite{rungepacket,oldframework}.  We do not repeat that extension.
Here we give a different certificate, adapted to equations with several
polynomial branches.  It applies directly to mixed polynomials in two
variables and supplies both coordinate bounds, including the exceptional
fibres where a remainder polynomial vanishes.

The arithmetic principle is elementary: a nonzero integer factor of an
integer \(w\) has absolute value at most \(|w|\).  If two factors differ
by a polynomial that eventually grows faster than \(w\), an integral
solution cannot lie beyond an explicit cutoff.  This gives a producer
for a restricted, syntactically recognizable class of height bounds.
It makes no completeness claim for certificate search and does not supply
the missing general Siegel or Faltings height theorem
\cite{rungepacket,faltingspacket}.

\subsection{The finite certificate and its estimates}

For an integer polynomial \(Q(X)=\sum_{k=0}^s q_kX^k\), write
\[
 \|Q\|_1=\sum_{k=0}^s|q_k|,
 \qquad
 E_Q(H)=\sum_{k=0}^s |q_k|H^k.
\]
Trailing zero coefficients are removed before computing a degree.
The zero polynomial has the empty coefficient array; every degree
comparison below is made only after checking that the relevant
polynomials are nonzero.

\begin{definition}[A certificate from separated polynomial branches]
\label{def:split-cert}
The finite input consists of integer polynomials
\(A_0,\ldots,A_{m-1},R\), with \(m\geq2\) and \(R\ne0\), and indices
\(i,j<m\).  Form \(P=A_i-A_j\), and require
\[
 d=\deg P>r=\deg R.
\]
In particular \(P\ne0\), \(d\geq1\), and the chosen branches are
different.  Put
\begin{align}
 L&=\sum_{k<d}|p_k|,
 & C&=\|R\|_1,\label{eq:split-LC}\\
 H&=\max\{1,L+2C\},
 & K&=\max_{t<m}E_{A_t}(H)+E_R(H).
 \label{eq:split-HK}
\end{align}
All these conditions and numbers are computed from the coefficient
arrays.  Since \(R\ne0\), the integer \(C\) is at least one.
\end{definition}

\begin{lemma}[Polynomial estimates]
\label{lem:split-poly-estimates}
For \(|x|\leq H\), with \(H\geq1\), one has
\( |Q(x)|\leq E_Q(H)\).
If \(Q\) has degree \(s\geq1\), nonzero integer leading coefficient,
and \(M=|x|\geq1\), then
\begin{equation}
 M^s\leq |Q(x)|+
          \left(\sum_{k<s}|q_k|\right)M^{s-1}.
 \label{eq:split-dominance}
\end{equation}
Consequently a nonzero integer polynomial \(Q\) has no integer root
with \(|x|>\|Q\|_1\).
\end{lemma}

\begin{proof}
The triangle inequality applied to the finite evaluation sum gives
\[
 |Q(x)|\leq\sum_{k=0}^s|q_k||x|^k
              \leq\sum_{k=0}^s|q_k|H^k.
\]
For the second assertion, isolate the leading term.  Because
\(|q_s|\geq1\),
\[
 M^s\leq |q_sx^s|
 \leq |Q(x)|+\sum_{k<s}|q_k|M^k
 \leq |Q(x)|+
          \left(\sum_{k<s}|q_k|\right)M^{s-1}.
\]
The final comparison uses \(M\geq1\).
If \(Q(x)=0\), cancellation of the positive integer \(M^{s-1}\)
yields \(M\leq\sum_{k<s}|q_k|\leq\|Q\|_1\).
A nonzero constant polynomial has no root at all, so it supplies the
remaining case.  All cancellations concern positive integers; no real
root, asymptotic notation, or logarithm is needed.
\end{proof}

\begin{lemma}[Size of an integer factor]
\label{lem:split-factor}
If \(u,v,w\in\Z\), \(uv=w\), and \(w\ne0\), then
\(u\ne0\), \(v\ne0\), and \(|u|\leq|w|\).
If a finite product of integers is zero, some factor is zero.
\end{lemma}

\begin{proof}
Either zero factor would force \(w=0\).  A nonzero integer \(v\)
satisfies \(|v|\geq1\), so
\(|w|=|u||v|\geq|u|\).  For the last assertion, induct through the
finite product using the integer product-zero law.  This also yields a
bounded algorithm: scan the factor list for its first zero.  Integer
equality is decidable, so this finite search does not use an infinitary
choice principle.
\end{proof}

\begin{theorem}[Explicit bounds from separated branches]
\label{thm:split-bound}
For every certificate in Definition~\ref{def:split-cert}, every integer
solution of
\begin{equation}
 \prod_{t<m}\bigl(y-A_t(x)\bigr)=R(x)
 \label{eq:split-equation}
\end{equation}
satisfies \(|x|\leq H\) and \(|y|\leq K\), with the explicit
bounds in \eqref{eq:split-HK}.
There is a uniform primitive recursive function of the certificate
that returns a complete finite list of its integer solutions.
\end{theorem}

\begin{proof}
First prove the bound for \(x\).  Integer order is decidable; in the
case \(|x|\leq H\) there is nothing to do.  In the other case set
\(M=|x|>H\).  Since \(H\geq2C\), we have \(M>C\).
Lemma~\ref{lem:split-poly-estimates} applied to the nonzero
polynomial \(R\) shows that \(R(x)\ne0\).

Put \(u_t=y-A_t(x)\).  For a fixed \(t<m\), the product of all
other factors is an integer \(v_t\), and
\(u_tv_t=R(x)\ne0\).  Lemma~\ref{lem:split-factor} therefore gives
\( |u_t|\leq|R(x)|\).  Applying this to the two selected indices,
\begin{equation}
 |P(x)|=|u_j-u_i|
       \leq |u_j|+|u_i|\leq2|R(x)|.
 \label{eq:split-factor-upper}
\end{equation}
The polynomial estimates and \(r<d\) imply
\[
 M^d\leq |P(x)|+LM^{d-1},
 \qquad
 |R(x)|\leq CM^r\leq CM^{d-1}.
\]
Combining these with \eqref{eq:split-factor-upper},
\[
 M^d\leq (L+2C)M^{d-1}.
\]
The factor \(M^{d-1}\) is positive, so cancellation gives
\(M\leq L+2C\leq H\), contradicting the case assumption.
Thus \(|x|\leq H\).

For the \(y\)-bound, the cases \(R(x)\ne0\) and \(R(x)=0\)
are both necessary.  In the first case, the factor estimate applies
without any lower bound on \(|x|\), and in particular
\[
 |y-A_0(x)|\leq |R(x)|.
\]
Using the just-proved \(x\)-bound,
\[
 |y|\leq|A_0(x)|+|R(x)|
       \leq\max_{t<m}E_{A_t}(H)+E_R(H)=K.
\]
In the second case the finite product in
\eqref{eq:split-equation} is zero.
Lemma~\ref{lem:split-factor} gives an index \(t<m\) with
\(y=A_t(x)\).  Therefore
\( |y|\leq E_{A_t}(H)\leq K\).
This includes all exceptional fibres over roots of \(R\).

Finally, enumerate \(x=-H,-H+1,\ldots,H\) and, for each such \(x\),
enumerate \(y=-K,-K+1,\ldots,K\).  Retain a pair exactly when the
two integers in \eqref{eq:split-equation} are equal.  The two loop
lengths are \(2H+1\) and \(2K+1\), and their product is an explicit
bound on the number of tests.  Every returned pair is a solution by
its acceptance test.  Every solution occurs because of the two bounds
already proved.  Polynomial evaluation, the bounds, finite products,
and these bounded loops are primitive recursive uniformly in their
finite coefficient inputs.
\end{proof}

\begin{remark}[How to attach the certificate to another equation]
Suppose an integer polynomial \(F(X,Y)\) is given together with a
nonzero integer \(D\) and a coefficientwise checked identity
\[
 D F(X,Y)=\prod_{t<m}\bigl(Y-A_t(X)\bigr)-R(X).
\]
Then its integer zero set is exactly the set in
\eqref{eq:split-equation}: one implication substitutes \(F=0\),
and the converse cancels the nonzero integer \(D\).
Thus a coefficient identity connects the certificate to the intended
input equation.  Verifying the identity is a finite computation;
proving that every equation of interest admits such an identity would
be a separate theorem, and is not asserted here.
\end{remark}

For two branches this method is closely related to completing a
square.  With three or more branches, mixed powers of \(Y\) generally
remain, and the proof works directly with them.  This is an extension
of the certificate classes treated explicitly in the inspected packets,
not a claim of historical novelty for Runge's method.  It also has a
different purpose from the existing Mordell--Weil descent theorem,
which already produces generators and coordinates under actual
quotient-cover and contraction hypotheses \cite{mordellpacket}.
Neither input interface generates the missing hypotheses of the other.

\subsection{A complete family of mixed equations}

The preceding theorem gives a uniform height bound in a family whose
integral solutions can in fact be classified exactly.

\begin{theorem}[An exact classification for every number of branches]
\label{thm:split-classification}
For every integer \(m\geq3\), the complete integer solution set of
\begin{equation}
 \prod_{j=1}^{m}\bigl(y-x^j\bigr)=1
 \label{eq:split-family}
\end{equation}
consists of the following points:
\begin{enumerate}
 \item \((0,1)\) and \((1,2)\), for every \(m\);
 \item \((0,-1)\) and \((1,0)\), in addition, precisely when
       \(m\) is even;
 \item \((-1,0)\), in addition, precisely when
       \(m\bmod4\) is \(0\) or \(1\).
\end{enumerate}
The list is a primitive recursive function of \(m\).
\end{theorem}

\begin{proof}
Every integer factor of a product equal to \(1\) is \(1\) or \(-1\).
Indeed, its absolute value is at most one by
Lemma~\ref{lem:split-factor}, and it cannot be zero.
Since \(m\geq3\), the first and third factors are present, so
\[
 |x^3-x|
 =|(y-x)-(y-x^3)|\leq2.
\]
For \(x\geq2\), however,
\(x^3-x=x(x-1)(x+1)\geq2\cdot1\cdot3=6\).
For \(x\leq-2\), apply the same inequality to \(-x\) and use the
oddness of \(X^3-X\).  Hence \(x\in\{-1,0,1\}\).

When \(x=0\), the equation is \(y^m=1\).
Its integer solutions are \(y=1\), and also \(y=-1\) for even
\(m\): a number of absolute value at least two has a power of
absolute value at least \(2^m>1\), and zero has power zero.
When \(x=1\), exactly the same argument applied to
\((y-1)^m=1\) gives \(y=2\), and also \(y=0\) for even \(m\).

When \(x=-1\), the first two factors are \(y+1\) and \(y-1\).
They must both belong to \(\{-1,1\}\).  Since their difference is
two, the only possibility is \(y+1=1\), \(y-1=-1\), hence \(y=0\).
At this point every odd-indexed factor is \(1\), and every even-indexed
factor is \(-1\).  The product is consequently
\[
 (-1)^{\lfloor m/2\rfloor}.
\]
Write \(m=4q+s\), where \(0\leq s<4\).  Then
\(\lfloor m/2\rfloor=2q+\lfloor s/2\rfloor\), which is even
exactly for \(s=0,1\).  This proves the stated condition.

The same substitutions verify every listed point.  Thus the list is
sound and complete, and computing it needs only parity and remainder
modulo four.  The proof consists of finite integer inequalities and
bounded product identities, uniformly in \(m\).
\end{proof}

At \(m=3\), the equation is explicitly
\[
 y^3-(x+x^2+x^3)y^2+(x^3+x^4+x^5)y-x^6=1,
\]
and its only integer solutions are \((0,1)\) and \((1,2)\).
The general certificate uses the branches \(X\) and \(X^3\),
so that \(L=C=1\); it gives \(H=3\) and \(K=3^m+1\).
The direct argument in Theorem~\ref{thm:split-classification} sharpens
that coarse box to the exact list.  No assertion of smoothness,
irreducibility, or genus is used in this arithmetic classification.

\subsection{Primitive recursion and the induction audit}
\label{subsec:split-audit}

The certificate is numerical data for one fixed algorithm, rather than
a code for an arbitrary program to be evaluated.  Choose canonical
signed-integer codes and finite arrays.  Polynomial evaluation uses
Horner recursion.  Subtraction pads with zero; multiplication uses bounded
convolution; an \(m\)-fold product is a primitive recursion in \(m\).
Trailing-zero removal is a bounded scan.  Absolute coefficient sums,
weighted sums, finite maxima, and the rectangular solution scan are
bounded folds.  Thus the implementation is genuinely uniform in degree,
coefficient size, and the number of branches.

The proofs require the following quantifier-free invariants after
expanding their fixed primitive recursive symbols:
\begin{enumerate}
 \item A Horner prefix or suffix equals its corresponding finite
       evaluation sum; addition and convolution respect evaluation.
 \item A processed product has the asserted factorization.  Nonzero
       total product entails that every processed factor is nonzero,
       and every product of nonzero integers has absolute value at
       least one.
 \item Partial coefficient sums dominate partial absolute evaluations;
       for \(M\geq1\), a smaller exponent gives a no-larger power.
 \item The scan for a zero factor either stores a verified zero index
       or has checked every preceding factor to be nonzero.
 \item The solution-list accumulator contains exactly the successful
       pairs in the already processed part of the explicit box.
 \item At \((-1,0)\), the product through index \(j\) equals
       \((-1)^{\lfloor j/2\rfloor}\).
\end{enumerate}
Bounded statements about earlier indices are represented by Boolean
folds.  Their use does not enlarge the induction schema.
The large degree, coefficient, or enumeration bounds cause no
proof-theoretic difficulty: primitive recursion places no polynomial-time
restriction on an algorithm.

The supplied \texttt{SplitTargets.lean} defines named pending goals for
the factor bound, polynomial evaluation bound, dominance estimate, large-root
exclusion, the two height bounds, complete enumeration and uniform family
classification. Their definitions require integer coefficient arrays, normalized degree,
Horner evaluation, branch products, the Boolean certificate checker,
and the two bounds in \eqref{eq:split-HK}.  The checker must compute
\(A_i-A_j\) itself; trusting separately supplied difference coefficients
would leave a gap.

The displayed arguments give a concrete route to \(\ipra\): establish
the signed-integer and array equations, prove the local update lemmas,
apply quantifier-free induction to the stated invariants, and then use
intuitionistic case analysis on the decidable arithmetic guards.
Producing the corresponding \texttt{Prf PRAAx} objects remains part of
the formalization work.  Native Lean statements, Python evaluations, and
coefficient checks alone do not discharge that object-language obligation.

The accompanying reference program uses integer arithmetic only.  It
checks malformed certificates, expands the mixed cubic independently,
enumerates several complete certificate boxes, includes examples where
\(R(x)=0\), and compares the family classification for \(m=3,\ldots,12\)
on explicit small boxes.  These recorded checks validate implementation
behavior.  The global completeness claims rest on the proofs above,
not on the extent of a numerical search.

\section{Formalization status and a concrete completion plan}
\label{sec:formalization}

The report contains complete written proofs of its stated restricted
mathematical results, with an explicit arithmetic induction audit. The
archive adds executable definitions, exact finite checks and a checked
part of the embedded proof infrastructure. Those categories have different
meanings; the following inventory makes each claim reviewable.

\subsection{What was checked}

The standalone Lean files were compiled with Lean 4.32.0, using only the
standard distribution. The compiler commit was
\begin{center}\small\code{8c9756b28d64dab099da31a4c09229a9e6a2ef35}.\end{center} The arithmetic kernel is a source-preserving extraction from
the pinned repository, with imports adapted as described in
Section~\ref{sec:object-iteration}. The build includes:
\begin{longtable}{@{}>{\raggedright\arraybackslash}p{.27\linewidth}>{\raggedright\arraybackslash}p{.35\linewidth}>{\raggedright\arraybackslash}p{.32\linewidth}@{}}
\toprule
File & Content & Verification meaning \\
\midrule\endhead
\code{ObjectKernel.lean} & Unchanged syntax, natural deduction and iPRA
declarations from three upstream modules & Kernel checked after import
adaptation; not new mathematical results \\
\code{ObjectIteration.lean} & Substitution lemmas and the general
primitive-recursion invariant rule & Actual embedded \code{Prf PRAAx}
proof transformation \\
\code{ObjectExample.lean} & Iterated successor equals its iteration count &
Closed embedded arithmetic derivation using the new rule \\
\code{NegativeSyntax.lean} & Negative translation, syntax laws, preservation
of QF, translation of all source axioms & Proved syntax facts and actual
embedded derivations of translated axioms; full proof translation pending \\
\code{ODETargets.lean} & Polynomial tables, Euler algorithm, solution name,
coefficient certificate and exact goal propositions & Definitions and target
types checked; the goal propositions are unproved \\
\code{SplitTargets.lean} & Integer polynomial algorithms, certificate,
bounds, enumeration and goal propositions & Definitions and target types
checked; the goal propositions are unproved \\
\code{Audit.lean} & Axiom reports and explicit target-type checks & Records
dependencies; checking a target's type does not establish it \\
\bottomrule
\end{longtable}

The audited new theorems use Lean's standard \code{propext} and
\code{Quot.sound} in the metatheory; some require only the former. The
reports contain no \code{sorryAx} or \code{Classical.choice} dependency.
Those metatheoretic principles do not add excluded middle, arbitrary
induction or choice to the embedded \code{PRAAx} proof calculus. The
source theory \code{ClassicalPRAAx} is an explicitly defined datatype of
axioms used as the translation input, not a new trusted Lean axiom.

The full repository native adapter was inspected against the current API
but was not built against the whole repository dependency closure. It is
clearly separated from the standalone build. No full object-theory proof
of the ODE construction, separated-branch bound or logical proof translator
is claimed.

\subsection{Definitions supplied for the analytic development}

\code{ODETargets.lean} uses finite coordinate functions
\code{Fin d -> Rat} for rational vectors. A monomial consists of a rational
coefficient and a finite exponent vector; a polynomial is a list of
monomials. Repeated exponent vectors are allowed. The polynomial-zero
checker adds the coefficients belonging to each exponent vector occurring
in the list and checks that all those sums vanish. Thus it handles
cancellation, and does not use equality of functions on an infinite
domain as a computational test.

The file defines the exact Euler-node recursion, rational polygon
interpolation, coefficient magnitude and Lipschitz bounds, the automatic
local horizon, the integer growth and Cauchy constants, and the reindexed
name. The predicate \code{Controlled} is precisely the finite-precision
solution contract in Definition~\ref{def:controlled-flow}.
\code{SameName} states agreement at arbitrarily accurate common
precisions. Its use is metatheoretic: each fixed evaluator becomes a PR
function symbol in the embedded version. The main target names are
\code{RefinementGoal}, \code{LocalExistenceGoal}, \code{UniquenessGoal}
and \code{ConservationGoal}.

The targets quantify over finite polynomial input data, and the uniform
algorithm evaluates that one finite grammar. This accommodates variable
dimension and degree without an unclocked interpreter for arbitrary PR
syntax.

\subsection{Dependency order for the main proofs}

Here is a path with independently reviewable endpoints, rather than a list
of unproved high-level axioms.

\paragraph{Stage 1: numerical representations.}
Implement natural encodings for integers, positive-denominator rationals,
finite arrays and polynomial tables. Prove the local ring and order laws,
finite-product absolute-value rules and positive-denominator cancellation
inside \code{Prf PRAAx}. Implement Boolean folds for bounded assertions.
The invariants in Appendix~\ref{sec:coding} describe the numerical
induction formulas. The existing repository may supply parts of this
library, but their actual proof endpoints must be imported explicitly.

\paragraph{Stage 2: reusable finite estimates.}
For ODEs, formalize coefficient bounds, rational Bernoulli, the finite
binomial estimate $(1+1/N)^N\le3$, and the variable-error Gronwall fold.
Each is an algebraic inequality with bounded index conditions. For split
branches, formalize \code{intFactorAbsLeGoal},
\code{polyEvalAbsLeGoal}, \code{polyDominanceGoal} and
\code{nonzeroPolyNoLargeRootGoal}. These are concrete named propositions
in \code{SplitTargets.lean}; their mathematical proofs are in
Section~\ref{sec:split-branches}.

The general rule of Theorem~\ref{thm:object-rule} discharges the outer
primitive-recursion induction after these local step proofs are supplied.
It is especially useful for folding bounded maxima, products, Euler tables
and accumulated errors. It does not replace the arithmetic in the local
step proofs.

\paragraph{Stage 3: the ODE invariant and refinement.}
Formalize the node bound
$\norm{x_{n,j}-x_0}\le jhM$, with the guard $j\le2^n$ encoded as a
Boolean or a quantifier-free implication. The paired-mesh state stores
both current nodes, parity of the fine-node index, and accumulated
error. Prove the recurrence in Theorem~\ref{thm:ode-refine}, apply the
finite Gronwall bound, and then extend from endpoints by affine
interpolation. A second finite fold telescopes refinements.
This proves \code{RefinementGoal} and the regularity component of
\code{LocalExistenceGoal}.

\paragraph{Stage 4: ODE meaning and conservation.}
The local-defect proof uses a finite partition of a rational interval by
mesh points, exact midpoint identities, and a rational triangle inequality.
It completes the controlled-solution contract. For uniqueness, fix the
requested accuracy first, compute $N,p$, and compare the finite table of
two candidates at that precision. For conservation, prove coefficient
checker soundness and the multivariate binomial remainder, then use one
accumulated-error invariant. Finally introduce the accuracy challenge
universally. No induction formula asserts equality of arbitrary real
functions.

\paragraph{Stage 5: complete integral-point lists.}
Prove \code{splitXBoundGoal} by coefficient dominance, and
\code{splitYBoundGoal} by the two cases $R(x)=0$ and $R(x)\ne0$.
Signed-interval membership and a rectangle scan give
\code{enumerateSolutionsIffGoal}. The numerical witnesses for a point in
the rectangle are $x+H$ and $y+K$; their bounds are already known.
The general family classification then needs only the integer unit
factors, $x\in\{-1,0,1\}$, and the parity fold described in its proof.

\paragraph{Stage 6: logical translation and native transport.}
The syntax and axiom layers of negative translation are already proved.
The next new lemma is the arithmetic equality-stability adapter, followed
by formula stability and induction on a complete source derivation.
The proof cases are fully spelled out in Section~\ref{sec:logic}.
After producing the desired \code{Prf PRAAx} theorem, apply the existing
\code{pra\_iff\_native} bridge. Its use is the final step, not an
alternative to Stages 1--5.

\subsection{Finite reference computations}

The exact-rational ODE program checks 5,342 assertions over exponential,
oscillator, nonlinear local-flow, Hermite-accumulator and circle-area
examples. It checks the finite inequalities and four coefficientwise
conservation certificates, and rejects an invalid certificate. The
nonlinear example uses small meshes because exact rational numerators and
denominators grow rapidly; the proof itself covers arbitrary precision.
No floating-point estimate is used as a proof oracle.

The integer program checks the complete certificate rectangle for the
mixed cubic, several further accepted certificates, exceptional fibers
where the remainder vanishes, coefficient expansion consistency, the
explicit family at $m=3,\ldots,12$ on stated finite boxes, and malformed
certificate rejection. The uniform classification and the absence of
solutions outside the certificate rectangle follow from the written
theorems, not from finite testing.

The source files and evidence record how to reproduce these runs. These
programs are reference implementations; neither is an iPRA proof checker.

\subsection{The difficulties that remain}

The framework makes the missing work more specific. It does not infer a
general height bound from finiteness or a cardinality bound, manufacture
Mordell--Weil quotient representatives, or establish the missing uniform
analytic estimates in prime-distribution packets. It does not decide
whether an arbitrary represented real is zero, choose singular Puiseux
branches from arbitrary Cauchy inputs, or infer a Cauchy modulus from finite
window stability. The source reports' warnings about these operations
remain applicable.

What is resolved here is narrower and concrete: an actual embedded
iteration rule, a written logical transfer mechanism with a checked axiom
layer, a constructed polynomial-flow interface for a stated analytic
unification problem, and a further explicit class of integral-point
problems with complete primitive-recursive solution lists.

\appendix
\section{Numerical coding and the induction discipline}\label{sec:coding}

This appendix fixes enough representation detail to make the route from
the mathematical arguments to first-order iPRA unambiguous. The particular
codes can be changed by proved PR translations. Their logical role cannot
be omitted.

\subsection{Total numerical operations and validity guards}

Use the signed coding
\[
 \operatorname{enc}(z)=
 \begin{cases}2z,&z\ge0,\\2|z|-1,&z<0.\end{cases}
\]
Decoding, sign, absolute value, comparison and the integer ring operations
are PR. Their laws reduce to natural addition, multiplication
and truncated subtraction. A rational is a pair $(z,d)$ with $d>0$;
canonical reduction by the gcd is convenient but unnecessary for the
arithmetic estimates. Equality and order use cross multiplication with
positive denominators. On a malformed rational code, use a fixed default
value and keep the validity flag explicit. Division in the ODE algorithm
always has a separately proved positive denominator.

For a concrete pairing, one can use
$\langle a,b\rangle=2^a(2b+1)-1$. Bounded trial division and a bounded
search recover its components. The decoder is total. A finite sequence
can be coded by a length together with nested pairs. Access to entry $i$
is a recursion bounded by the length, returning a default past the end.
Alternatively, a table whose entries are bounded by $C$ can be coded in
base $C+1$. Then a length-$r$ table has code below $(C+1)^r$.
Both encodings and the usual append, map, fold and bounded search
operations are PR.

Canonical polynomial codes remove trailing zero coefficients. In the
split-branch checker, the remainder must be nonzero, and the chosen
difference is recomputed from the two selected input polynomials.
Comparison of normalized lengths expresses the strict degree gap even
for constant remainders. The zero polynomial is not assigned a genuine
degree and then accidentally admitted as a certificate.

\subsection{Why variable size does not introduce a universal evaluator}

The interpreter for a finite integer or rational polynomial table is one
fixed PR program. It performs bounded folds over coefficients, monomials
and coordinates, using numerical exponentiation. It does not evaluate
arbitrary PR program descriptions. Iterating this fixed update a variable
number of times is precisely the primitive recursion operation allowed
in the language.

For example, let $s_{j+1}=\operatorname{Step}(s_j,c)$ where $c$ is a
polynomial-code parameter. The function $(j,c,s_0)\mapsto s_j$ is PR
because \code{Step} is one fixed PR function, even if its outputs become
very large. Its evaluation time need not be elementary or practical.
Nothing here asserts that the evaluator of all PR syntax is itself PR.

The same distinction applies to Lean's function-valued notation
\code{Fin d -> Rat}. This is a convenient finite coordinate presentation
in the metatheory. The object theory uses its explicit finite sequence
code, and the correspondence is proved by entrywise decoding lemmas.
Arbitrary functions $\N\to\Q$ are never inserted into the numerical
state as if they were finite data.

\subsection{Bounded formulas and finite induction}

For a fixed PR Boolean test $b(\vec x,i)$, define
\[
 A(\vec x,0)=1,\qquad
 A(\vec x,j+1)=A(\vec x,j)\mathbin{\wedge_b}b(\vec x,j),
\]
where $\wedge_b$ is conjunction on $0,1$. A finite induction
proves that $A(\vec x,j)=1$ expresses success at every index $i<j$.
For the forward direction, fix an arbitrary challenge $i$ and induct on
$j$ with the quantifier-free invariant
$(i<j\land A(\vec x,j)=1)\to b(\vec x,i)=1$.
For the reverse direction, fix a terminal bound $N$ and keep
$\forall i<N\,b(\vec x,i)=1$ as an assumption independent of the
induction variable. Induct on $j$ with
$j\le N\to A(\vec x,j)=1$. Universal elimination from the fixed
assumption supplies the successor test. Thus even the proof of bounded
reflection does not presuppose induction on a quantified formula.
The corresponding disjunctive fold and first-success register implement
bounded existence and witness selection. With atomic equality and order
tests, structural induction on a fixed bounded formula proves its
equivalence to a Boolean equation. This is the reflection interface used
throughout; it is not reflection for all theorems of iPRA.

If a bounded loop has guard $j\le N$, its invariant can have the form
\[
 (j\le N\longrightarrow B(s_j,j,\vec x)=1).
\]
Here order is a PR test. A stopped transition freezes the state outside
the active range. The induction formula is quantifier-free. In particular,
the assertion that the first $j$ entries all satisfy a bound must be
expressed by its Boolean prefix fold before invoking \code{PRA.qfInd}.

The principal invariants in the new applications are:
\begin{longtable}{@{}>{\raggedright\arraybackslash}p{.22\linewidth}>{\raggedright\arraybackslash}p{.42\linewidth}>{\raggedright\arraybackslash}p{.30\linewidth}@{}}
\toprule
Loop & Numerical invariant & Necessary guard or auxiliary data \\
\midrule\endhead
Polynomial evaluation & Horner prefix agrees with the corresponding finite
sum & Encoded coefficient length and prefix index \\
Polynomial bound & Absolute partial sum is bounded by the sum of absolute
terms & Nonnegative power base; leading coefficient nonzero when used \\
Euler trajectory & $\norm{x_j-x_0}\le jhM$ & $j\le N$,
$hN=T$, and the box bound \\
Error propagation & $e_j\le r^j(e_0+\sum_{i<j}\delta_i)$ & $r\ge1$
and nonnegative rational errors \\
Conservation & Accumulated discrepancy is at most the sum
of quadratic step remainders & Coefficientwise Lie-derivative identity \\
Split factors & Each selected nonzero integer factor has absolute value
at most the product & Full product nonzero; zero-product case separate \\
Rectangle scan & Output contains exactly the successful tests already
visited & Explicit bounds $H,K$ and scan position \\
\bottomrule
\end{longtable}

\subsection{Accuracy challenges and rational gap elimination}

An equality of names is an assertion about all requested precisions. The
correct proof order is to fix an arbitrary precision $k$, compute a finite
schedule from $k$, prove the corresponding rational inequality by the
preceding induction discipline, and then universally introduce $k$.
One does not induct on equality of represented values.

Removing a vanishing numerical error is elementary. Suppose
rational $a,b,C$ with $C\ge0$ satisfy $a\le b+C2^{-n}$ for all $n$.
If $a>b$, write the positive rational gap $a-b=u/v$ with $u,v\ge1$.
Choose an integer $J\ge Cv/u$ and take $n=J+1$. The finite inequality
$2^{J+1}>J$ makes $C2^{-n}<u/v$ when $C>0$, a contradiction.
For $C=0$ the conclusion is immediate. The choice of $J$ is an explicit
rational ceiling, and the exponent bound is PR. This removes a rational
slack; it does not decide the sign of an arbitrary represented real.

Finally, every theorem described as uniform has one finite numerical
algorithm and one arithmetic proof schema for that algorithm. A collection
of unrelated standard-numeral proofs is not substituted for a single
variable-parameter proof.

\begin{thebibliography}{99}
\small
\bibitem{snapshot} All repository references below are pinned to commit
\path{d56cd5505fb93135b90b048bf3dcca9e73615162}.
The accompanying \path{provenance/source_manifest.json} gives exact paths,
Git blob identifiers and SHA-256 hashes for the consulted sources.
\bibitem{repository} ProveIt, Primitive recursive arithmetic: scope and status.
\href{https://github.com/VladimirReshetnikov/ProveIt/blob/d56cd5505fb93135b90b048bf3dcca9e73615162/Logic/PrimitiveRecursiveArithmetic/README.md}{Pinned repository source}.
\bibitem{kernel} ProveIt, ConciseIPRA arithmetic kernel: PR syntax and axioms.
\href{https://github.com/VladimirReshetnikov/ProveIt/blob/d56cd5505fb93135b90b048bf3dcca9e73615162/Logic/PrimitiveRecursiveArithmetic/ConciseIPRA/Lean/ConciseIPRA/PR.lean}{Pinned repository source}.
\bibitem{bridge} ProveIt, Native bi-interpretation: pra\_iff\_native.
\href{https://github.com/VladimirReshetnikov/ProveIt/blob/d56cd5505fb93135b90b048bf3dcca9e73615162/Logic/PrimitiveRecursiveArithmetic/ConciseIPRA/Lean/ConciseIPRA/NativeBiInterpretation.lean}{Pinned repository source}.
\bibitem{oldframework} ProveIt, Finite contracts for mathematics in intuitionistic primitive recursive arithmetic.
\href{https://github.com/VladimirReshetnikov/ProveIt/blob/d56cd5505fb93135b90b048bf3dcca9e73615162/Logic/PrimitiveRecursiveArithmetic/docs/proofs/foundations/ipra_framework/README.md}{Pinned repository source}.
\bibitem{embeddingpacket} ProveIt, Mathematics in iPRA through finite observations.
\href{https://github.com/VladimirReshetnikov/ProveIt/blob/d56cd5505fb93135b90b048bf3dcca9e73615162/Logic/PrimitiveRecursiveArithmetic/docs/proofs/foundations/ipra_embedding_framework/README.md}{Pinned repository source}.
\bibitem{finiteobservation} ProveIt, Finite observations and uniform constructions in iPRA.
\href{https://github.com/VladimirReshetnikov/ProveIt/blob/d56cd5505fb93135b90b048bf3dcca9e73615162/Logic/PrimitiveRecursiveArithmetic/docs/proofs/foundations/ipra_finite_observation/README.txt}{Pinned repository source}.
\bibitem{evidencepacket} ProveIt, Evidence-indexed finite mathematics in intuitionistic primitive recursive arithmetic.
\href{https://github.com/VladimirReshetnikov/ProveIt/blob/d56cd5505fb93135b90b048bf3dcca9e73615162/Logic/PrimitiveRecursiveArithmetic/docs/proofs/foundations/ipra_evidence_framework/report.tex}{Pinned repository source}.
\bibitem{uniformpacket} ProveIt, Uniform certificate framework.
\href{https://github.com/VladimirReshetnikov/ProveIt/blob/d56cd5505fb93135b90b048bf3dcca9e73615162/Logic/PrimitiveRecursiveArithmetic/docs/proofs/foundations/ipra_uniform_certificates/framework.tex}{Pinned repository source}.
\bibitem{frameworkdated} ProveIt, iPRA certificate-and-modulus framework.
\href{https://github.com/VladimirReshetnikov/ProveIt/blob/d56cd5505fb93135b90b048bf3dcca9e73615162/Logic/PrimitiveRecursiveArithmetic/docs/proofs/foundations/ipra_framework_2026-11-15/README.md}{Pinned repository source}.
\bibitem{impossibilitiespacket} ProveIt, Three classical impossibilities in iPRA; especially Further research questions, item 2.
\href{https://github.com/VladimirReshetnikov/ProveIt/blob/d56cd5505fb93135b90b048bf3dcca9e73615162/Logic/PrimitiveRecursiveArithmetic/docs/proofs/algebra/three_impossibilities_ipra/three_classical_impossibilities_ipra.tex}{Pinned repository source}.
\bibitem{rungepacket} ProveIt, Siegel packet: a finite Runge certificate and the general height-bound boundary.
\href{https://github.com/VladimirReshetnikov/ProveIt/blob/d56cd5505fb93135b90b048bf3dcca9e73615162/Logic/PrimitiveRecursiveArithmetic/docs/proofs/diophantine-number-theory/siegel_ipra/paper/03_runge.tex}{Pinned repository source}.
\bibitem{mordellpacket} ProveIt, A primitive recursive Mordell--Weil descent schema.
\href{https://github.com/VladimirReshetnikov/ProveIt/blob/d56cd5505fb93135b90b048bf3dcca9e73615162/Logic/PrimitiveRecursiveArithmetic/docs/proofs/diophantine-number-theory/mordell_weil_ipra/paper/02_descent.tex}{Pinned repository source}.
\bibitem{faltingspacket} ProveIt, Faltings theorem and intuitionistic PRA: status and missing positive bound.
\href{https://github.com/VladimirReshetnikov/ProveIt/blob/d56cd5505fb93135b90b048bf3dcca9e73615162/Logic/PrimitiveRecursiveArithmetic/docs/proofs/diophantine-number-theory/faltings_ipra/README.md}{Pinned repository source}.
\bibitem{moschovakis} Joan Rand Moschovakis,
\emph{Calibrating the negative interpretation}, arXiv:2101.10313v4, 2021.
\url{https://arxiv.org/abs/2101.10313v4}.
Used for the background negative-translation method; the restricted PRA
argument in Section~\ref{sec:logic} is proved explicitly in this report.
\end{thebibliography}

\end{document}
